%%%%%%%%%%  scan idx 372  |  folio 359  %%%%%%%%%%

Ces
relations étant isobares, le quotient est une \(K\)-algèbre graduée \(\Omega\). On définit un morphisme \(\Lambda\xrightarrow{\widetilde c_0}\Ch(\Omega)\) par :
\[
  x\rightsquigarrow
  \bigl(\epsilon(x),1+c^1(x-\epsilon(x))+\cdots\bigr).
\]
Les relations imposées montrent que c'est une théorie de Chern de \(\Lambda\) à valeurs dans \(\Omega\). Si de plus on a une théorie de Chern \(f\), à valeurs dans une \(K\)-algèbre graduée \(A\), on a de façon unique un homomorphisme : \(\Omega\to A\), tel que \(f\) se factorise par \(\Ch(\Omega)\); il est défini en envoyant, pour tout \(i\ge1\), et tout \(x\in\Fil^1\Lambda\), l'élément \(c^i(x)\) sur le \(i\)-ième coefficient de \(f(x)\) (ce qui est bien défini car \(f\) est un \(\lambda\)-homomorphisme).

Supposons maintenant \(K\) de caractéristique nulle, et \(\Lambda\) de filtration discrète~; soit \(n\) tel que \(\Fil^n\Lambda=0\). Alors l'algèbre graduée \(\Omega\) est nulle en degrés \(\ge n\). Pour le montrer, il faut montrer que pour \(x_1,\ldots,x_p\in\Fil^1\Lambda\), et \(k_1+\cdots+k_p\ge n\), on a :
\[
  c^{k_1}(x_1)\cdots c^{k_p}(x_p)=0.
\]
Or par hypothèse, on a :
\[
  \gamma^{k_1}(x_1)\cdots\gamma^{k_p}(x_p)=0\ ;
\]
En appliquant le \(\lambda\)-homomorphisme \(\widetilde c_0\), on obtient :
\[
  \gamma^{k_1}(\widetilde c_0(x_1))\cdots
  \gamma^{k_p}(\widetilde c_0(x_p))=(0,1).
\]
Or d'après (6.6.4) et (6.5.3), ceci entraine :
\[
  (k_1+\cdots+k_p)!\,
  c^{k_1}(x_1)\cdots c^{k_p}(x_p)=0.
\]
Comme \(K\) est de caractéristique nulle, on a le résultat voulu.

La propriété universelle de \(\Omega\) appliqué à l'homomorphisme canonique \(\widetilde c\) donne un diagramme commutatif :
